\leavevmode\printedpage{64}
% head: A. Grothendieck : Dérivateurs, Ch. VII, Édition des Universités de Montpellier II et Paris VII
\msp{page 99}

\section*{10\quad Comparison question for \textmd{Cat}-fibered categories satisfying the ``homotopy theorem'' (filtering base)}

We are now going to generalize to the case of a (Cat)-fibered category over another
\[
\begin{tikzcd}
C \arrow[d, "\pi"'] \\
T,
\end{tikzcd}
\]
thinking mainly of the case $T = \uTy$, $C = \uCh_{\Sigma}(X;a,b)$. The essential hypothesis is the following:

\noindent\begin{tabular}{@{}p{0.3\textwidth}p{0.64\textwidth}@{}}
(Homotopy hypothesis) & If $t,t' \in \Ob T$ are such that $\Hom(t,t') \neq \emptyset$, then the functors $\varphi^{*} : C_{t'} \longrightarrow C_{t}$ associated to the $\varphi \in \Hom(t,t')$ are homotopic.
\end{tabular}

They therefore define one and the same arrow in Cat\,hot. I denote by $T_{\mathrm{ord}}$ or $\To$ the set $\Ob(T)$, with as order relation the one deduced from the category structure:
\[
t \leq t' \qquad\Longleftrightarrow\qquad \Hom_{T}(t,t') \neq \emptyset,
\]
so that we have a functor
\[
T \longrightarrow \To
\]
which is \textit{universal} for functors from $T$ into ``ordered'' categories (i.e. where every arrow is monomorphic \texttt{[no, cf. }$\bullet\rightrightarrows\bullet$\texttt{; rather, where always }$\#\Hom(a,b) \leq 1$ \texttt{(admitting nonidentity isomorphisms)]}) or for functors $F$ on $T$ such that $F$ is constant on every $\Hom(t,t')$. The homotopy hypothesis can be expressed by saying that the composite functor
\[
\begin{tikzcd}
T^{o} \arrow[r] & \Cat \arrow[r] & \Cathot \\
t \arrow[r, mapsto] & C_{t} &
\end{tikzcd}
\]
\msp{page 100}
factors (obviously in a unique way) as a functor
\[
H_{C} : \To^{o} \longrightarrow \Cathot,
\]
i.e. that we have a commutative diagram
\printedpage{65}
% head: A. Grothendieck : Dérivateurs, Ch. VII, Édition des Universités de Montpellier II et Paris VII
\[
\begin{tikzcd}
T^{o} \arrow[r, "t\,\longmapsto\, C_{t}"] \arrow[d] & \Cat \arrow[d] \\
\To^{o} \arrow[r, "H_{C}"] & \Cathot.
\end{tikzcd}
\]

We are going to assume in addition:

\noindent\begin{tabular}{@{}p{0.36\textwidth}p{0.58\textwidth}@{}}
(Filtering hypothesis) & The ordered set $\To$ associated to $T$ is downward \textit{filtering} (left filtering), i.e. if $t,t' \in \Ob T$, there exist $t'' \in \Ob T$ and arrows $t'' \longrightarrow t$, $t'' \longrightarrow t'$.
\end{tabular}

This hypothesis is indeed satisfied in $\uTy$, since for $\tau,\tau' \in \Ob \uTy = \mathrm{Ty}$, it suffices to take $\tau\tau'$, which maps to $\tau$ by $\mathrm{id}_{\tau} * p_{\tau'}$, to $\tau'$ by $p_{\tau} * \mathrm{id}_{\tau'}$, where $p_{\tau} : \tau \longrightarrow 1$ is the morphism from $\tau$ to the final object.

Thus, $H_{C}$ defines an ind-object of Cat\,hot, which we are going to denote by $\underline{\underline{C}}$. If we assume in addition:

\noindent\begin{tabular}{@{}p{0.36\textwidth}p{0.58\textwidth}@{}}
(Countability hypothesis, always satisfied if $\mathrm{card}\,\To \leq \aleph_{0}$) & $\To$ has a \textit{countable} subset cofinal on the left.
\end{tabular}

then this ind-object of Cat\,hot can be defined, in many ways, in terms
\msp{page 101}
of a direct system in Cat (not just in Cat\,hot)
\[
(C_{n})_{n\in\mathbf{N}},
\]
from which we recover the ind-object $\underline{\underline{C}}$ in Cat\,hot by passing to the corresponding direct system in Cat\,hot, then to the associated ind-object.

To construct the direct system $C_{n}$, it suffices to take an infinite chain
\[
t_{0} \longleftarrow t_{1} \longleftarrow t_{2} \longleftarrow \;\cdots\; \longleftarrow t_{n} \longleftarrow t_{n+1} \longleftarrow \;\cdots
\]
of morphisms in $T$, such that the set of the $t_{i}$ is cofinal on the left in $\To$. We are going to assume in addition that $C$ is \textit{strictly fibered} over $T$ (i.e. coming from a genuine functor $T^{o} \longrightarrow \Cat$, i.e. we can choose the inverse image functors
\[
\varphi^{*} : C_{t'} \longrightarrow C_{t}
\]
in a strictly transitive way, not just transitive up to isomorphism) - a case to which we can always reduce, moreover, by the trick of \textsc{Giraud}. This hypothesis is moreover verified in the case $\uCh_{\Sigma}(X;a,b)$ over $\uTy$. Thus, the
\[
(C_{t_{i}})_{i\in\mathbf{N}}
\]
form a direct system of categories indexed by $\mathbf{N}$, the transition morphism
\[
C_{n} \longrightarrow C_{n'} \qquad \text{i.e. } C_{t_{n}} \longrightarrow C_{t_{n'}}
\]
\printedpage{66}
% head: A. Grothendieck : Dérivateurs, Ch. VII, Édition des Universités de Montpellier II et Paris VII
\msp{page 102}
being the inverse image functor $\varphi^{*}_{n,n'}$, where
\[
\varphi_{n,n'} : t_{n'} \longrightarrow t_{n}
\]
is the composite of the $t_{n'} \longrightarrow t_{n'-1} \longrightarrow \;\cdots\; \longrightarrow t_{n}$. (If we did not have strict transitivity, we would define $\varphi^{*}_{n,n'}$ as the composite of the inverse image functors along successive transition arrows $t_{n'} \longrightarrow t_{n'-1} \longrightarrow \;\cdots\; \longrightarrow t_{n+1} \longrightarrow t_{n}$, this would be $\varphi^{*}_{n,n'}$ only up to canonical isomorphism.)

Let
\[
C_{\infty} = \varinjlim_{n} C_{n} = \varinjlim_{n} C_{t_{n}}
\]
be the category $\varinjlim$ . (It will be especially nice if we have been able to choose the system $(t_{n})_{n\in\mathbf{N}}$ in such a way that the transition functors
\[
C_{t_{n}} \longrightarrow C_{t_{n-1}}
\]
are fully faithful, and perhaps even \texttt{[some]} monomorphisms, in such a way that the $C_{n}$ identify with full subcategories of their direct limit $C_{\infty}$.) I would like to understand the precise relations between the following three kinds of objects:

\begin{itemize}[leftmargin=2.5em]
\item[\circled{a}] $C$ itself (as an object in Cat, not of $\Cat/T$), and especially the objects of Cat\,hot and of Hot which correspond to it.

\item[\circled{b}] \msp{page 103}$C_{\infty}$, and the objects of Cat\,hot and of Hot which correspond to it. To what extent are these latter % [sic] deniers -> derniers, missing letter
independent % [sic] indépendant -> indépendants, agreement with the plural subject
of the choice of the inverse system cofinal on the left $(t_{n})_{n\in\mathbf{N}}$ in $T$? What is the role of $C_{\infty}$? Is it canonically isomorphic to $C$ in Cat\,hot, or in Hot? (In which case the question of its ``role'' would seem to me resolved!)

\item[\circled{c}] The ind-object % [sic] Le ind-objet -> L'ind-objet, missing elision
$H_{C}$ in Cat\,hot, and the ind-object of Hot which it defines.

I note that the ind-object $H_{C}$, and also the categories $C_{\infty}$ constructed as above, do not change when we replace $T$ by a subcategory $T'$ cofinal on the left, or also by a category of the form $T/t_{0}$ ($t_{0} \in \Ob T$), which gives rise to the associated ordered set % [sic] ordonnée -> ordonné, agreement with "ensemble"
\[
(T/t)_{\mathrm{o}} = \To/t = \{t \in \To \mid t \leq t_{0}\}
\]
\texttt{[rather} $(T/t_{0})_{\mathrm{o}} = \To/t_{0} = \{t \in \To \mid t \leq t_{0}\}$\texttt{]}. If the relation between the objects in a), and those in b) and in c), is very close, we would conclude from this % [sic] concluerait -> conclurait, spelling of the conditional
that the ``type'' of $C$ in Cat\,hot, or in Hot as the case may be, does not change on replacing $C$ by
\printedpage{67}
% head: A. Grothendieck : Dérivateurs, Ch. VII, Édition des Universités de Montpellier II et Paris VII
\msp{page 104}
$C \times_{T} T'$ or by $C/t_{0} \simeq C \times_{T} T/t_{0}$. Better still, it would not change on replacing $T$ by $(\mathbf{N})^{o}$, the (ordered) category opposite to $\mathbf{N}$ (thus having a unique final object), and the fibered category $C$ over $T$ by the category $C_{\underline{t}}$ induced by $C$ over $(\mathbf{N})^{o}$, via the functor
\[
\underline{t} : (\mathbf{N})^{o} \longrightarrow T
\]
defined by the inverse system $\underline{t} = (t_{n})_{n\in\mathbf{N}}$ in $T$. This would thus mean that to understand the homotopy type (in the strict sense, i.e. in Cat\,hot, or broad, i.e. in Hot) of $C$, we can ``clean'' it in a draconian way, the base $T$ being viewed as a sort of somewhat muddy gangue, or at least plethoric.

\item[\circled{d}] In any case, if $T'$ is another left filtering category, and
\[
\alpha : T' \longrightarrow T
\]
a functor cofinal on the left, so that $C$ over $T$, and $C' = C \times_{T} T'$ over $T'$, define essentially the ``same'' object of Ind\,Cat\,hot, and the category $C'_{\infty}$ for $C'$ over $T'$ is among the categories $C_{\infty}$ for $C$ over $T$ - we may ask
\msp{page 105}
whether the canonical functor
\[
C' \;\longrightarrow\; C
\]
is also always an isomorphism in Cat\,hot, or if not, at least in Hot.
\end{itemize}

A particular interest attaches to $\pi_{0}(C)$, and moreover to its variance for $C$ over $T$ varying in the following way. If $S$ is in Cat, we look at the categories
\[
\uHom_{\Cat}(S,C_{t}) \qquad (t \in T),
\]
we find a functor
\[
t \longmapsto \uHom_{\Cat}(S,C_{t}), \qquad T^{o} \longrightarrow (\Cat),
\]
whence a fibered category over $T$, which I am going to denote by
\[
\sHom_{T}(S_{T},C)
\]
(Unless I am mistaken, it is interpreted as the fibered category over $T$ corresponding to the stack on $\mathrm{Top}(T) = T^{\wedge}$ of the morphism of $T$-stacks from $S_{T}$ to $C$, where $S_{T} = S \times T$ (considered as a fibered category over $T$ with essentially constant fiber $S$) to $C$. Here I identify fibered category over $T$, and stacks on the topos $T^{\wedge}$.) It is immediate that the homotopy property (p. 99) is stable under passage from $C$ to $\sHom_{T}(S_{T},C)$. We are interested in the variation of
\[
\pi_{0}(\sHom_{T}(S_{T},C))
\]
\printedpage{68}
% head: A. Grothendieck : Dérivateurs, Ch. VII, Édition des Universités de Montpellier II et Paris VII
\msp{page 106}
as a contrafunctor in $S$,
\[
H : S \longmapsto \pi_{0}\sHom_{T}(S_{T},C) \qquad (\Cat)^{o} \longrightarrow (\mathrm{Ens}),
\]
and in a possible expression of this contrafunctor $H$ in terms of the ind-object $H_{C}$ of Cat\,hot, by the hypothetical formula
\[
H(S) \simeq \Hom_{\mathrm{Ind}(\Cathot)}(\gamma(S),H_{C}),
\]
where
\[
\gamma : \Cat \longrightarrow \Cathot
\]
is the canonical functor. If we had such a canonical isomorphism functorial in $S$, this would give a description (up to unique isomorphism) of the ind-object $H_{C}$ in terms of the preceding functor $H$ on Cat.

A doubt comes to me here, in view of the example $T = \uTy$, $C = \uCh_{\Sigma}(X;a,b)$. The three types of objects a) b) c) have the same $\pi_{0}$, and it is certain that ``it's the right one''. On the other hand, for a variable parameter category $S$, it is also certain that the right functor in $S$ is
\[
S \longmapsto \pi_{0}(\uHom_{T}(S_{T},C)) \simeq \Hom_{\mathrm{Ind}\,\Cathot}(S,H_{C}),
\]
also isomorphic, \textit{if $S$ is finite,} to
\[
\Hom_{\Cathot}(S,C_{\infty}),
\]
which shows that, at least for arguments $S \in (\Cat)$ which are \textit{finite} categories, the objects
\msp{page 107}
$C_{\infty}$ and $H_{C}$ are ``good'', they allow us % [sic] il nous permettent -> ils nous permettent, agreement with the plural subject
to describe an important functor which dominates the situation under consideration. Under a very strong hypothesis, but one satisfied in the important cases, we can in these two formulas replace Cat\,hot by Hot.

By contrast, except for the case where $S$ is reduced to a point and where $\Hom_{\Cathot}$ or $\Hom_{\Hot}$ from $S$ to $C$ is none other than $\pi_{0}(C)$ ($\simeq \pi_{0}(C_{\infty}) \simeq \pi_{0}(H_{C})$), thus indeed ``good'', I do not know whether
\[
\underbrace{\Hom_{\Cathot}(S,C)}_{=\,\pi_{0}\uHom(S,C)} \qquad \text{or} \qquad \Hom_{\Hot}(S,C)
\]
is related in a reasonable way to the preceding sets (determined by $H_{C}$, resp. by $C_{\infty}$), e.g. whether they are isomorphic to it in a natural way, and whether they have a reasonable geometrico-topological meaning. I am beginning to suspect \textit{not} - hence that we would have here a case where a relation of a very general type concerning a $\pi_{0}$, here thus
\[
\pi_{0}(C) \simeq \pi_{0}(C_{\infty}) \simeq \pi_{0}(H_{C}),
\]
does not extend to an isomorphism on the higher $\pi_{i}$, and to a property of Hot-equivalence.